\section{Mechanism 2 in the gate world: a kept state that imagination cannot use}
\label{sec:planning}

The gate world is built so that mechanism 1 is off the table. The background is black, the wall is
large and slow, and the only uncertainty is which opening will be revealed. A model that keeps the
state must show the approaching wall; a model that averages over the two openings should still show
the shared solid part. Figure~\ref{fig:gate-filmstrip} shows what the grid JEPA does with a pre-reveal
context: decoded encodings of the \emph{real} future keep the wall and its opening at every step,
and the predictor's own rollout deletes the wall from the first predicted step, in all three
encoder seeds (Appendix~\ref{app:historical-gate}). The flow model's sampled rollout keeps the wall
and commits to one opening. This is the visible signature of the second mechanism: the state was
there, and the mean-trained predictor did not carry it forward.

\begin{figure}[t]
\centering
\paperfig{gate_s1_filmstrip.png}{1.0}
\caption{\textbf{The wall survives encoding and dies in prediction.} From one pre-reveal context
and the recorded actions: the simulator; the JEPA's decoded encodings of the real future; the
JEPA's decoded rollout; the flow model's decoded encodings and one sampled future. JEPA model
seed 1, the encoder used in the behavioral comparison below; flow seed 0; first eligible held-out
episode, chosen without looking at model outputs. A decoding failure on predicted latents mixes
prediction error with off-distribution read-out; the behavioral comparison is what turns this
picture into a claim.}
\label{fig:gate-filmstrip}
\end{figure}

\paragraph{Why a shared pixel policy.}
Earlier attempts to compare the two models through latent-input policies were uninformative in a
specific way: a latent-MLP policy trained by PPO fails the gate even on the \emph{exact} simulator
with true rewards (0 of 45 policies above $50\%$ survival; Appendix~\ref{app:matched}), so its
failure inside imagination says nothing about the model. The comparison in this section therefore
uses one pixel CNN and GRU policy for both families, fed by each model's own decoder, with the same
PPO, 150 updates of 512 rollouts of 48 steps, three policy seeds, final checkpoints, and 256
evaluation episodes on shared keys (Section~\ref{sec:worlds}). The JEPA encoder is model seed 1,
the seed on which decoded-latent control had succeeded during development; the flow model is one
dynamics checkpoint. Both facts limit the claim to these fitted models. The JEPA's reward and
termination heads are refitted on all 500,000 training labels so that head supervision is
count-matched to the flow recipe.

\paragraph{Interface qualification.}
Before imagination is read, the same policy is trained on real transitions seen through each
decoder. \pendingresult{Decoded-real controls: JEPA seeds 0/1 at 0.977/0.980 survival so far; flow
seeds 0/1 at 0.957/0.516. Insert all three seeds per family with walls passed, and state which
interface qualifies. If flow's decoded interface is unstable across seeds, say so here and carry it
into the imagination reading.}

\paragraph{Imagination.}
\pendingresult{Insert the paired JEPA-versus-flow imagination outcomes (three policy seeds each):
survival, walls passed, control length. Read them against the qualification row. Do not attribute a
gap specifically to conditional-mean averaging; the recipes differ in supervision stage.}

\paragraph{What the historical record adds.}
Across the earlier, unmatched protocols, none of 27 policies trained inside the JEPA's imagination
discovered the gate strategy, including nine that used a sampling head fitted on the predictor's
residual and nine that averaged eight of its samples, while policies trained on real transitions
through the same frozen encoders did (Table~\ref{tab:gate-main}, Appendix~\ref{app:historical-gate}).
The flow reference's historical success was a single policy of six. Both facts are why this section
rests on the matched comparison above rather than on those tables. An exact nine-branch calculation
on frozen encoders shows the scoring operator of Equation~\ref{eq:variance-cost} changing the
optimal action for $1$--$3\%$ of context--goal pairs (Appendix~\ref{app:oracle-variance}); it is a
small direct illustration, not a control result.
